\section{Scope: crops, groups and countries}\label{sec:scope}

\subsection{Crops and groups}
The index covers nine crops in three predefined groups:
\begin{itemize}
\item \textbf{cereals}: wheat, rice, maize;
\item \textbf{pulses}: chickpea, lentil, common bean;
\item \textbf{oilseeds}: soybean, groundnut, rapeseed-mustard.
\end{itemize}
The groups were fixed before any result was computed. A group score is the simple mean of its three crops'
scores (Section~\ref{sec:aggregation}). It describes these three crops, not every cereal, pulse or oilseed.

Each crop must be identified in five datasets: FAOSTAT production, the SPAM crop map, the Sacks calendar, the FAO
food balance sheets and the food composition table. The identifiers are listed in Table~\ref{tab:cropmap}
(Section~\ref{sec:data}). Three definitions need explanation.
\begin{itemize}
\item \textbf{Rapeseed-mustard} is a single item definition used everywhere: FAOSTAT ``Rape or colza seed'' (India
reports its rapeseed-mustard here; the separate ``Mustard seed'' item is about 1\% of the total and is excluded),
SPAM \texttt{RAPE}, and the food-balance items ``Rape and Mustardseed'' and ``Rape and Mustard Oil''.
\item \textbf{Common bean} is FAOSTAT ``Beans, dry''. In India and Myanmar this item also contains tropical
\emph{Vigna} pulses, which are not \emph{Phaseolus vulgaris}. The bean heat parameter is therefore applied beyond
the species it was measured on (Appendix~\ref{app:ledger}).
\item \textbf{Chickpea and lentil} have production and crop-map entries of their own. The food balance sheets,
however, report them only inside the aggregate ``Pulses, Other and products''. Their food supply is taken from the
FAOSTAT Supply Utilization Accounts (Section~\ref{sec:fbsmap}).
\end{itemize}

\subsection{Country selection rule}\label{sec:panelrule}
Countries are chosen per crop by the rule presented at the mid-evaluation, applied unchanged. Production is the
FAOSTAT mean for 2020--2024.
\begin{enumerate}
\item The \textbf{core countries} USA, Brazil, India and China are included for every crop they produce.
\item \textbf{One European and one African country.} A country is used for all crops if it is among its
continent's top three producers of every crop. Otherwise each crop uses its continent's largest producer.
\item The crop's \textbf{world top-3 producers} are added.
\end{enumerate}
A core country enters a crop's panel only if FAOSTAT records production there in 2020--2024 \emph{and} SPAM~2010
maps crop area there. With nine crops, no European or African country is in its continent's top three for every
crop, so the rule's own fallback applies. This changes two of the original four-crop panels: rice uses Italy
instead of Russia, and maize uses Ukraine instead of Russia. Three combinations are excluded by the rule:
\begin{itemize}
\item Brazil for chickpea and for lentil (no production or no mapped area);
\item groundnut in Europe (no European production recorded).
\end{itemize}
The result is \nPairs\ crop--country pairs in \nCountries\ countries (Table~\ref{tab:panels}). The panels cover
between half and nine-tenths of world production of each crop. A panel is a sample of the main producing
systems, not a census, so a crop's score describes these countries only.

\begin{table}[H]\centering\small
\caption{Country panels (FAOSTAT production 2020--24; world rank order; $*$ = less than 0.5\% of world output).}\label{tab:panels}
\adjustbox{max width=\linewidth}{\begin{tabular}{lp{10.5cm}r}
\toprule
Crop & Countries (world rank order; * = <0.5\% of world output) & \% world \\
\midrule
Wheat & China, India, Russia, USA, Egypt, Brazil & 51 \\
Rice & China, India, Bangladesh, Brazil, USA, Nigeria, Italy* & 63 \\
Maize & USA, China, Brazil, India, Ukraine, South Africa & 70 \\
Chickpea & India, Australia, Turkey, Ethiopia, Russia, USA, China* & 89 \\
Lentil & Canada, India, Australia, USA, Russia, China, Ethiopia & 82 \\
Common bean & India, Brazil, Myanmar, Tanzania, China, USA, Belarus & 59 \\
Soybean & Brazil, USA, Argentina, China, India, Russia, South Africa & 90 \\
Groundnut & China, India, Nigeria, USA, Brazil & 67 \\
Rapeseed-mustard & Canada, China, India, France, USA, South Africa*, Brazil* & 61 \\
\bottomrule
\end{tabular}
}
\end{table}

\subsection{Periods}
\begin{itemize}
\item \textbf{Present:} 2001--2024, observed monthly climate.
\item \textbf{Historical baseline:} 1971--2000, observed.
\item \textbf{Future:} 2041--2060, one climate model (MRI-ESM2-0) under one scenario (SSP2-4.5).
\end{itemize}
Importance is fixed at the 2021--2023 food balance sheets in every period. Crop area, irrigation and calendars are
fixed at their 2010 maps. Sensitivity is fixed at its literature values.
